% --- Archon physics formalization source begin ---
% archon:physics
% archon:covers PhyXMiniProblems/problem_phyx_mini_0633.lean
% archon:source-report reports/phyx_mini/problem_phyx_mini_0633.source.json

\chapter{Physics problem phyx\_mini\_0633}
\label{ch:PhyXMiniProblems_problem_phyx_mini_0633}

\paragraph{Problem source.}
\#\# Physical scenario

Peggy is standing in the center of a long, flat railroad car that has firecrackers tied to both ends. The car moves past Ryan, who is standing on the ground, with velocity \$v = 0.8c\$. Flashes from the exploding firecrackers reach him simultaneously \$1.0\textbackslash{};\textbackslash{}mu s\$ after the instant that Peggy passes him, and he later finds burn marks on the track \$300\textbackslash{};m\$ to either side of where he had been standing.

\#\# Question

According to Peggy, what times do the explosions occur relative to the time that Ryan passes her?

\#\# Answer choices

- A: A: \$-1.13\textbackslash{};\textbackslash{}mu s\$

- B: B: \$-1.33\textbackslash{};\textbackslash{}mu s\$

- C: C: \$-1.53\textbackslash{};\textbackslash{}mu s\$

- D: D: \$-1.73\textbackslash{};\textbackslash{}mu s\$

\#\# Figure caption (auxiliary; use the image as primary evidence)

The image illustrates a scenario involving two reference frames, Frame \textbackslash{}(S\textbackslash{}) and Frame \textbackslash{}(S'\textbackslash{}). 

- **Frame \textbackslash{}(S\textbackslash{}):** 
  - Associated with a person named Ryan.
  - Ryan is standing on the ground, looking at the train platform.

- **Frame \textbackslash{}(S'\textbackslash{}):**
  - Associated with a person named Peggy.
  - Peggy is standing on a platform (representing a train car) that is moving to the right with velocity \textbackslash{}(v\textbackslash{}).

- **Events:**
  - **Event L:** Marked with a red starburst on the left side of the platform.
    - Coordinates: \textbackslash{}((x\_L, t\_L) = (-300 \textbackslash{}, \textbackslash{}text\{m\}, 0 \textbackslash{}, \textbackslash{}text\{s\})\textbackslash{}).

  - **Event R:** Marked with a red starburst on the right side of the platform.
    - Coordinates: \textbackslash{}((x\_R, t\_R) = (300 \textbackslash{}, \textbackslash{}text\{m\}, 0 \textbackslash{}, \textbackslash{}text\{s\})\textbackslash{}).

- **Text:**
  - Blue text pointing to Peggy: "Peggy passes Ryan at \textbackslash{}(t = t' = 0\textbackslash{})."

- **Visual Elements:**
  - Horizontal line marked with positions: \textbackslash{}(-300 \textbackslash{}, \textbackslash{}text\{m\}\textbackslash{}), \textbackslash{}(0\textbackslash{}), and \textbackslash{}(300 \textbackslash{}, \textbackslash{}text\{m\}\textbackslash{}).
  - Ryan and Peggy are positioned to show relative motion between the reference frames.
  - An arrow labeled \textbackslash{}(v\textbackslash{}) indicates the direction of motion of the platform to the right.
  - Railway tracks are depicted under the platform.

This diagram likely demonstrates concepts of special relativity, illustrating how events are seen from different reference frames.

\#\# Dataset metadata

Relativity; Physical Model Grounding Reasoning, Spatial Relation Reasoning

\paragraph{Recorded answer/context.}
B: B: \$-1.33\textbackslash{};\textbackslash{}mu s\$

\paragraph{Figure/image path.}
/root/proposal\_for\_physic/science-mango/phyx\_mini\_run/phyx\_data/test\_image/633.png

\paragraph{Formalization target.}
create a compiling Lean file with sorry bodies at `PhyXMiniProblems/problem\_phyx\_mini\_0633.lean`. The Lean declarations must preserve the physical quantities, dimensions or dimensional roles, figure labels, governing-law hypotheses, and final relation expressed by this problem.
Use Mathlib/Physlib names found through LeanExplore where available. If a physics API is missing, introduce faithful local abstractions rather than scalar placeholder aliases.

\begin{theorem}[Physics formalization target]
\label{thm:physics:phyx_mini_0633:target}
\lean{PhyXMiniProblems.ProblemPhyXMini0633.problem_phyx_mini_0633}
\uses{def:physics:phyx-mini-0633:phyxminiproblems-problemphyxmini0633-railroadcarrelativitysetup, def:physics:phyx-mini-0633:phyxminiproblems-problemphyxmini0633-peggyexplosiontimemicroseconds, def:physics:phyx-mini-0633:phyxminiproblems-problemphyxmini0633-usesmetermicrosecondcoordinates, def:physics:phyx-mini-0633:phyxminiproblems-problemphyxmini0633-matchessuppliedrailroadcarfigure, def:physics:phyx-mini-0633:phyxminiproblems-problemphyxmini0633-matchesproblemreadouts, def:physics:phyx-mini-0633:phyxminiproblems-problemphyxmini0633-hasphysicalrelativisticparameters, def:physics:phyx-mini-0633:phyxminiproblems-problemphyxmini0633-satisfiesvacuumlightpropagation, def:physics:phyx-mini-0633:phyxminiproblems-problemphyxmini0633-satisfiesryantopeggylorentzboost, def:physics:phyx-mini-0633:phyxminiproblems-problemphyxmini0633-predictedpeggyexplosiontimemagnitudemicroseconds, def:physics:phyx-mini-0633:phyxminiproblems-problemphyxmini0633-recordeddatasetanswer, def:physics:phyx-mini-0633:phyxminiproblems-problemphyxmini0633-isuniquematchinganswerchoice}
The assigned autoformalize agent should translate this physics problem into Lean declarations in the covered file, with theorem and lemma proof bodies written as `by sorry`.
\end{theorem}
\begin{proof}
This is an autoformalization task, not a proof task. Produce statements that can later be proved by the physics prover without weakening the physical model.
\end{proof}

% --- Archon declaration topology begin ---
\section{Declaration topology}
\begin{definition}[speed In Meters Per Second]
  \label{def:physics:phyx-mini-0633:phyxminiproblems-problemphyxmini0633-speedinmeterspersecond}
  \lean{PhyXMiniProblems.ProblemPhyXMini0633.speedInMetersPerSecond}
  Read a nonnegative physical speed in coherent SI units, metres per second
\end{definition}
\begin{proof}
  This is the declared model definition of speed In Meters Per Second.
\end{proof}
\begin{definition}[vacuum Speed Of Light In Meters Per Second]
  \label{def:physics:phyx-mini-0633:phyxminiproblems-problemphyxmini0633-vacuumspeedoflightinmeterspersecond}
  \lean{PhyXMiniProblems.ProblemPhyXMini0633.vacuumSpeedOfLightInMetersPerSecond}
  Physlib's exact vacuum speed of light, read in metres per second
\end{definition}
\begin{proof}
  This is the declared model definition of vacuum Speed Of Light In Meters Per Second.
\end{proof}
\begin{definition}[vacuum Speed Of Light In Meters Per Microsecond]
  \label{def:physics:phyx-mini-0633:phyxminiproblems-problemphyxmini0633-vacuumspeedoflightinmeterspermicrosecond}
  \lean{PhyXMiniProblems.ProblemPhyXMini0633.vacuumSpeedOfLightInMetersPerMicrosecond}
  \uses{def:physics:phyx-mini-0633:phyxminiproblems-problemphyxmini0633-vacuumspeedoflightinmeterspersecond}
  The exact vacuum speed of light, read in metres per microsecond
\end{definition}
\begin{proof}
  This is the declared model definition of vacuum Speed Of Light In Meters Per Microsecond.
\end{proof}
\begin{definition}[Inertial Frame Label]
  \label{def:physics:phyx-mini-0633:phyxminiproblems-problemphyxmini0633-inertialframelabel}
  \lean{PhyXMiniProblems.ProblemPhyXMini0633.InertialFrameLabel}
  The two inertial frames printed in the supplied image
\end{definition}
\begin{proof}
  This is the declared model definition of Inertial Frame Label.
\end{proof}
\begin{definition}[Observer Label]
  \label{def:physics:phyx-mini-0633:phyxminiproblems-problemphyxmini0633-observerlabel}
  \lean{PhyXMiniProblems.ProblemPhyXMini0633.ObserverLabel}
  The two observers associated with the frames
\end{definition}
\begin{proof}
  This is the declared model definition of Observer Label.
\end{proof}
\begin{definition}[Explosion Side]
  \label{def:physics:phyx-mini-0633:phyxminiproblems-problemphyxmini0633-explosionside}
  \lean{PhyXMiniProblems.ProblemPhyXMini0633.ExplosionSide}
  The left and right firecrackers, explosion events, and light flashes
\end{definition}
\begin{proof}
  This is the declared model definition of Explosion Side.
\end{proof}
\begin{definition}[Event Label]
  \label{def:physics:phyx-mini-0633:phyxminiproblems-problemphyxmini0633-eventlabel}
  \lean{PhyXMiniProblems.ProblemPhyXMini0633.EventLabel}
  \uses{def:physics:phyx-mini-0633:phyxminiproblems-problemphyxmini0633-explosionside}
  Events needed to state the scenario and its light-propagation evidence
\end{definition}
\begin{proof}
  This is the declared model definition of Event Label.
\end{proof}
\begin{definition}[Observer Location]
  \label{def:physics:phyx-mini-0633:phyxminiproblems-problemphyxmini0633-observerlocation}
  \lean{PhyXMiniProblems.ProblemPhyXMini0633.ObserverLocation}
  Qualitative location of each observer in the scenario
\end{definition}
\begin{proof}
  This is the declared model definition of Observer Location.
\end{proof}
\begin{definition}[Horizontal Direction]
  \label{def:physics:phyx-mini-0633:phyxminiproblems-problemphyxmini0633-horizontaldirection}
  \lean{PhyXMiniProblems.ProblemPhyXMini0633.HorizontalDirection}
  Horizontal direction relative to the orientation of the image
\end{definition}
\begin{proof}
  This is the declared model definition of Horizontal Direction.
\end{proof}
\begin{definition}[Figure Feature]
  \label{def:physics:phyx-mini-0633:phyxminiproblems-problemphyxmini0633-figurefeature}
  \lean{PhyXMiniProblems.ProblemPhyXMini0633.FigureFeature}
  Named visual elements in image phyx\_data/test\_image/633.png
\end{definition}
\begin{proof}
  This is the declared model definition of Figure Feature.
\end{proof}
\begin{definition}[Railroad Car Relativity Figure]
  \label{def:physics:phyx-mini-0633:phyxminiproblems-problemphyxmini0633-railroadcarrelativityfigure}
  \lean{PhyXMiniProblems.ProblemPhyXMini0633.RailroadCarRelativityFigure}
  \uses{def:physics:phyx-mini-0633:phyxminiproblems-problemphyxmini0633-inertialframelabel, def:physics:phyx-mini-0633:phyxminiproblems-problemphyxmini0633-observerlabel, def:physics:phyx-mini-0633:phyxminiproblems-problemphyxmini0633-explosionside, def:physics:phyx-mini-0633:phyxminiproblems-problemphyxmini0633-observerlocation, def:physics:phyx-mini-0633:phyxminiproblems-problemphyxmini0633-horizontaldirection, def:physics:phyx-mini-0633:phyxminiproblems-problemphyxmini0633-figurefeature}
  Literal qualitative and numerical content of the primary image. The figure's event times are printed in seconds, while the working coordinates below use microseconds; their common displayed value is zero
\end{definition}
\begin{proof}
  This is the declared model definition of Railroad Car Relativity Figure.
\end{proof}
\begin{definition}[Railroad Car Relativity Setup]
  \label{def:physics:phyx-mini-0633:phyxminiproblems-problemphyxmini0633-railroadcarrelativitysetup}
  \lean{PhyXMiniProblems.ProblemPhyXMini0633.RailroadCarRelativitySetup}
  \uses{def:physics:phyx-mini-0633:phyxminiproblems-problemphyxmini0633-inertialframelabel, def:physics:phyx-mini-0633:phyxminiproblems-problemphyxmini0633-explosionside, def:physics:phyx-mini-0633:phyxminiproblems-problemphyxmini0633-eventlabel, def:physics:phyx-mini-0633:phyxminiproblems-problemphyxmini0633-horizontaldirection, def:physics:phyx-mini-0633:phyxminiproblems-problemphyxmini0633-railroadcarrelativityfigure}
  Independent physical quantities and spacetime coordinates for the problem. No Peggy-frame explosion time is defined from an answer choice
\end{definition}
\begin{proof}
  This is the declared model definition of Railroad Car Relativity Setup.
\end{proof}
\begin{definition}[speed Fraction Of Light]
  \label{def:physics:phyx-mini-0633:phyxminiproblems-problemphyxmini0633-speedfractionoflight}
  \lean{PhyXMiniProblems.ProblemPhyXMini0633.speedFractionOfLight}
  \uses{def:physics:phyx-mini-0633:phyxminiproblems-problemphyxmini0633-speedinmeterspersecond, def:physics:phyx-mini-0633:phyxminiproblems-problemphyxmini0633-vacuumspeedoflightinmeterspersecond, def:physics:phyx-mini-0633:phyxminiproblems-problemphyxmini0633-railroadcarrelativitysetup}
  The dimensionless speed parameter β = v / c
\end{definition}
\begin{proof}
  This is the declared model definition of speed Fraction Of Light.
\end{proof}
\begin{definition}[time Coordinate Readout]
  \label{def:physics:phyx-mini-0633:phyxminiproblems-problemphyxmini0633-timecoordinatereadout}
  \lean{PhyXMiniProblems.ProblemPhyXMini0633.timeCoordinateReadout}
  \uses{def:physics:phyx-mini-0633:phyxminiproblems-problemphyxmini0633-inertialframelabel, def:physics:phyx-mini-0633:phyxminiproblems-problemphyxmini0633-eventlabel, def:physics:phyx-mini-0633:phyxminiproblems-problemphyxmini0633-railroadcarrelativitysetup}
  Time-coordinate readout in the time unit selected by the setup
\end{definition}
\begin{proof}
  This is the declared model definition of time Coordinate Readout.
\end{proof}
\begin{definition}[x Coordinate Readout]
  \label{def:physics:phyx-mini-0633:phyxminiproblems-problemphyxmini0633-xcoordinatereadout}
  \lean{PhyXMiniProblems.ProblemPhyXMini0633.xCoordinateReadout}
  \uses{def:physics:phyx-mini-0633:phyxminiproblems-problemphyxmini0633-inertialframelabel, def:physics:phyx-mini-0633:phyxminiproblems-problemphyxmini0633-eventlabel, def:physics:phyx-mini-0633:phyxminiproblems-problemphyxmini0633-railroadcarrelativitysetup}
  Axial position-coordinate readout in the length unit selected by the setup
\end{definition}
\begin{proof}
  This is the declared model definition of x Coordinate Readout.
\end{proof}
\begin{definition}[peggy Explosion Time Microseconds]
  \label{def:physics:phyx-mini-0633:phyxminiproblems-problemphyxmini0633-peggyexplosiontimemicroseconds}
  \lean{PhyXMiniProblems.ProblemPhyXMini0633.peggyExplosionTimeMicroseconds}
  \uses{def:physics:phyx-mini-0633:phyxminiproblems-problemphyxmini0633-explosionside, def:physics:phyx-mini-0633:phyxminiproblems-problemphyxmini0633-railroadcarrelativitysetup, def:physics:phyx-mini-0633:phyxminiproblems-problemphyxmini0633-timecoordinatereadout}
  Peggy's microsecond time readout for one of the two explosion events
\end{definition}
\begin{proof}
  This is the declared model definition of peggy Explosion Time Microseconds.
\end{proof}
\begin{definition}[Uses Meter Microsecond Coordinates]
  \label{def:physics:phyx-mini-0633:phyxminiproblems-problemphyxmini0633-usesmetermicrosecondcoordinates}
  \lean{PhyXMiniProblems.ProblemPhyXMini0633.UsesMeterMicrosecondCoordinates}
  \uses{def:physics:phyx-mini-0633:phyxminiproblems-problemphyxmini0633-vacuumspeedoflightinmeterspermicrosecond, def:physics:phyx-mini-0633:phyxminiproblems-problemphyxmini0633-railroadcarrelativitysetup}
  The setup uses metre and microsecond coordinates with the calibrated c
\end{definition}
\begin{proof}
  This is the declared model definition of Uses Meter Microsecond Coordinates.
\end{proof}
\begin{definition}[Matches Supplied Railroad Car Figure]
  \label{def:physics:phyx-mini-0633:phyxminiproblems-problemphyxmini0633-matchessuppliedrailroadcarfigure}
  \lean{PhyXMiniProblems.ProblemPhyXMini0633.MatchesSuppliedRailroadCarFigure}
  \uses{def:physics:phyx-mini-0633:phyxminiproblems-problemphyxmini0633-railroadcarrelativitysetup, def:physics:phyx-mini-0633:phyxminiproblems-problemphyxmini0633-timecoordinatereadout, def:physics:phyx-mini-0633:phyxminiproblems-problemphyxmini0633-xcoordinatereadout}
  Primary-image evidence: Ryan is in ground frame S, Peggy is at the center of the railroad car in frame S', the car moves right, the two explosions are at (x\_L,t\_L)=(-300 m,0 s) and (x\_R,t\_R)=(300 m,0 s), and the passing event is the common origin t=t'=0
\end{definition}
\begin{proof}
  This is the declared model definition of Matches Supplied Railroad Car Figure.
\end{proof}
\begin{definition}[Agrees With One Decimal Microsecond Readout]
  \label{def:physics:phyx-mini-0633:phyxminiproblems-problemphyxmini0633-agreeswithonedecimalmicrosecondreadout}
  \lean{PhyXMiniProblems.ProblemPhyXMini0633.AgreesWithOneDecimalMicrosecondReadout}
  Agreement with a time printed to one decimal place in microseconds
\end{definition}
\begin{proof}
  This is the declared model definition of Agrees With One Decimal Microsecond Readout.
\end{proof}
\begin{definition}[Matches Problem Readouts]
  \label{def:physics:phyx-mini-0633:phyxminiproblems-problemphyxmini0633-matchesproblemreadouts}
  \lean{PhyXMiniProblems.ProblemPhyXMini0633.MatchesProblemReadouts}
  \uses{def:physics:phyx-mini-0633:phyxminiproblems-problemphyxmini0633-railroadcarrelativitysetup, def:physics:phyx-mini-0633:phyxminiproblems-problemphyxmini0633-speedfractionoflight, def:physics:phyx-mini-0633:phyxminiproblems-problemphyxmini0633-timecoordinatereadout, def:physics:phyx-mini-0633:phyxminiproblems-problemphyxmini0633-xcoordinatereadout, def:physics:phyx-mini-0633:phyxminiproblems-problemphyxmini0633-agreeswithonedecimalmicrosecondreadout}
  Prose data not supplied by the coordinate labels alone: v=0.8c, both flashes reach Ryan together about 1.0 μs after passing, and the two burn marks lie 300 m on opposite sides of Ryan's initial position
\end{definition}
\begin{proof}
  This is the declared model definition of Matches Problem Readouts.
\end{proof}
\begin{definition}[Has Physical Relativistic Parameters]
  \label{def:physics:phyx-mini-0633:phyxminiproblems-problemphyxmini0633-hasphysicalrelativisticparameters}
  \lean{PhyXMiniProblems.ProblemPhyXMini0633.HasPhysicalRelativisticParameters}
  \uses{def:physics:phyx-mini-0633:phyxminiproblems-problemphyxmini0633-railroadcarrelativitysetup, def:physics:phyx-mini-0633:phyxminiproblems-problemphyxmini0633-speedfractionoflight, def:physics:phyx-mini-0633:phyxminiproblems-problemphyxmini0633-timecoordinatereadout}
  Positivity and subluminality conditions selecting the physical branch
\end{definition}
\begin{proof}
  This is the declared model definition of Has Physical Relativistic Parameters.
\end{proof}
\begin{definition}[Satisfies Vacuum Light Propagation]
  \label{def:physics:phyx-mini-0633:phyxminiproblems-problemphyxmini0633-satisfiesvacuumlightpropagation}
  \lean{PhyXMiniProblems.ProblemPhyXMini0633.SatisfiesVacuumLightPropagation}
  \uses{def:physics:phyx-mini-0633:phyxminiproblems-problemphyxmini0633-railroadcarrelativitysetup, def:physics:phyx-mini-0633:phyxminiproblems-problemphyxmini0633-timecoordinatereadout, def:physics:phyx-mini-0633:phyxminiproblems-problemphyxmini0633-xcoordinatereadout}
  Vacuum light propagation between each explosion and its arrival at Ryan. In the selected metre/microsecond coordinates, a lightlike interval obeys |Δx| = c Δt. This is a generic governing law and does not contain either Peggy-frame answer time
\end{definition}
\begin{proof}
  This is the declared model definition of Satisfies Vacuum Light Propagation.
\end{proof}
\begin{definition}[Satisfies Ryan To Peggy Lorentz Boost]
  \label{def:physics:phyx-mini-0633:phyxminiproblems-problemphyxmini0633-satisfiesryantopeggylorentzboost}
  \lean{PhyXMiniProblems.ProblemPhyXMini0633.SatisfiesRyanToPeggyLorentzBoost}
  \uses{def:physics:phyx-mini-0633:phyxminiproblems-problemphyxmini0633-railroadcarrelativitysetup, def:physics:phyx-mini-0633:phyxminiproblems-problemphyxmini0633-speedfractionoflight}
  The passive Lorentz boost from Ryan's ground coordinates to Peggy's car coordinates. Physlib's boost acts on the complete SpaceTime 1 event, with the standard time component γ(β) (ct - βx). The law is quantified over all events and contains no solved explosion-time value or answer-choice label
\end{definition}
\begin{proof}
  This is the declared model definition of Satisfies Ryan To Peggy Lorentz Boost.
\end{proof}
\begin{definition}[predicted Peggy Explosion Time Magnitude Microseconds]
  \label{def:physics:phyx-mini-0633:phyxminiproblems-problemphyxmini0633-predictedpeggyexplosiontimemagnitudemicroseconds}
  \lean{PhyXMiniProblems.ProblemPhyXMini0633.predictedPeggyExplosionTimeMagnitudeMicroseconds}
  \uses{def:physics:phyx-mini-0633:phyxminiproblems-problemphyxmini0633-vacuumspeedoflightinmeterspermicrosecond}
  Magnitude predicted from β=0.8, |x|=300 m, and exact vacuum c
\end{definition}
\begin{proof}
  This is the declared model definition of predicted Peggy Explosion Time Magnitude Microseconds.
\end{proof}
\begin{lemma}[peggy explosion times exact]
  \label{lem:physics:phyx-mini-0633:phyxminiproblems-problemphyxmini0633-peggy-explosion-times-exact}
  \lean{PhyXMiniProblems.ProblemPhyXMini0633.peggy_explosion_times_exact}
  \uses{def:physics:phyx-mini-0633:phyxminiproblems-problemphyxmini0633-railroadcarrelativitysetup, def:physics:phyx-mini-0633:phyxminiproblems-problemphyxmini0633-peggyexplosiontimemicroseconds, def:physics:phyx-mini-0633:phyxminiproblems-problemphyxmini0633-usesmetermicrosecondcoordinates, def:physics:phyx-mini-0633:phyxminiproblems-problemphyxmini0633-matchessuppliedrailroadcarfigure, def:physics:phyx-mini-0633:phyxminiproblems-problemphyxmini0633-matchesproblemreadouts, def:physics:phyx-mini-0633:phyxminiproblems-problemphyxmini0633-satisfiesryantopeggylorentzboost, def:physics:phyx-mini-0633:phyxminiproblems-problemphyxmini0633-predictedpeggyexplosiontimemagnitudemicroseconds}
  The two Ryan-simultaneous explosions are not simultaneous for Peggy: the left event is later and the right event is earlier by equal magnitudes. This is a derived conclusion, not a governing-law field
\end{lemma}
\begin{proof}
  The conclusion follows from the governing relations and figure readouts named in the statement of peggy explosion times exact.
\end{proof}
\begin{definition}[Answer Choice]
  \label{def:physics:phyx-mini-0633:phyxminiproblems-problemphyxmini0633-answerchoice}
  \lean{PhyXMiniProblems.ProblemPhyXMini0633.AnswerChoice}
  Labels of the four microsecond-valued choices in the source problem
\end{definition}
\begin{proof}
  This is the declared model definition of Answer Choice.
\end{proof}
\begin{definition}[displayed Peggy Time Microseconds]
  \label{def:physics:phyx-mini-0633:phyxminiproblems-problemphyxmini0633-displayedpeggytimemicroseconds}
  \lean{PhyXMiniProblems.ProblemPhyXMini0633.displayedPeggyTimeMicroseconds}
  \uses{def:physics:phyx-mini-0633:phyxminiproblems-problemphyxmini0633-answerchoice}
  The negative microsecond time printed beside each answer label
\end{definition}
\begin{proof}
  This is the declared model definition of displayed Peggy Time Microseconds.
\end{proof}
\begin{definition}[recorded Dataset Answer]
  \label{def:physics:phyx-mini-0633:phyxminiproblems-problemphyxmini0633-recordeddatasetanswer}
  \lean{PhyXMiniProblems.ProblemPhyXMini0633.recordedDatasetAnswer}
  \uses{def:physics:phyx-mini-0633:phyxminiproblems-problemphyxmini0633-answerchoice}
  The answer label recorded by the source dataset; this is not a premise
\end{definition}
\begin{proof}
  This is the declared model definition of recorded Dataset Answer.
\end{proof}
\begin{definition}[Rounds To Nearest Hundredth Microsecond]
  \label{def:physics:phyx-mini-0633:phyxminiproblems-problemphyxmini0633-roundstonearesthundredthmicrosecond}
  \lean{PhyXMiniProblems.ProblemPhyXMini0633.RoundsToNearestHundredthMicrosecond}
  Agreement with a time displayed to the nearest hundredth of a microsecond
\end{definition}
\begin{proof}
  This is the declared model definition of Rounds To Nearest Hundredth Microsecond.
\end{proof}
\begin{definition}[Matches Answer Choice]
  \label{def:physics:phyx-mini-0633:phyxminiproblems-problemphyxmini0633-matchesanswerchoice}
  \lean{PhyXMiniProblems.ProblemPhyXMini0633.MatchesAnswerChoice}
  \uses{def:physics:phyx-mini-0633:phyxminiproblems-problemphyxmini0633-railroadcarrelativitysetup, def:physics:phyx-mini-0633:phyxminiproblems-problemphyxmini0633-peggyexplosiontimemicroseconds, def:physics:phyx-mini-0633:phyxminiproblems-problemphyxmini0633-answerchoice, def:physics:phyx-mini-0633:phyxminiproblems-problemphyxmini0633-displayedpeggytimemicroseconds, def:physics:phyx-mini-0633:phyxminiproblems-problemphyxmini0633-roundstonearesthundredthmicrosecond}
  The negative choices compare to the right-hand explosion, the explosion that occurs before the passing event in Peggy's frame
\end{definition}
\begin{proof}
  This is the declared model definition of Matches Answer Choice.
\end{proof}
\begin{definition}[Is Unique Matching Answer Choice]
  \label{def:physics:phyx-mini-0633:phyxminiproblems-problemphyxmini0633-isuniquematchinganswerchoice}
  \lean{PhyXMiniProblems.ProblemPhyXMini0633.IsUniqueMatchingAnswerChoice}
  \uses{def:physics:phyx-mini-0633:phyxminiproblems-problemphyxmini0633-railroadcarrelativitysetup, def:physics:phyx-mini-0633:phyxminiproblems-problemphyxmini0633-answerchoice, def:physics:phyx-mini-0633:phyxminiproblems-problemphyxmini0633-matchesanswerchoice}
  The selected label is the unique displayed choice matching Peggy's time
\end{definition}
\begin{proof}
  This is the declared model definition of Is Unique Matching Answer Choice.
\end{proof}
% --- Archon declaration topology end ---
% --- Archon physics formalization source end ---
